%==============================================================================
% FICHAMENTO CRÍTICO — Artificial Intelligence for Assessment of
% Facial Attractiveness
% Conforme NBR 6023 (referências) e NBR 10520 (citações)
% Capítulo 20 — Atratividade Facial, Odontologia Estética e Ortodontia
%==============================================================================
\section{Fichamento: \citeonline{wang2024facial}}

% --- 1. IDENTIFICAÇÃO ---
\subsection{Identificação}
\begin{tabular}{ll}
\textbf{Título:} & Artificial intelligence for assessment of facial
attractiveness: a deep learning approach using Vision Transformers
and geometric morphometrics on 3D stereophotogrammetry \\
\textbf{Autor(es):} & Wang, Lei; Chen, Wei; Liu, Yuhang; Tanaka, Yuki;
Kang, Dongwoo; Park, Min-Ji; Li, Qingyi; Silva, Rafael N. \\
\textbf{Periódico:} & Oral Surgery, Oral Medicine, Oral Pathology and Oral Radiology \\
\textbf{Ano:} & 2024 \\
\textbf{Volume:} & 138 \\
\textbf{Número:} & 6 \\
\textbf{Páginas:} & 546--565 \\
\textbf{DOI:} & \url{https://doi.org/10.1016/j.oooo.2024.01.023} \\
\textbf{Qualis:} & A2 (Odontologia) \\
\textbf{Tipo:} & Artigo original de desenvolvimento e validação de modelo \\
\end{tabular}

% --- 2. OBJETIVO ---
\subsection{Objetivo}
Desenvolver e validar um sistema automatizado de avaliação de
atratividade facial baseado em Vision Transformers (ViT-L/16)
combinados com análise de morfometria geométrica tridimensional
(dos pontos landmarks faciais), utilizando imagens de
estereofotogrametria 3D (3dMD). O estudo visa superar a
subjetividade inerente às escalas de avaliação estética facial
tradicionais --- que apresentam alta variabilidade inter e
intraobservador, dependência de painéis de avaliadores treinados
e baixa reprodutibilidade ---, estabelecendo métricas objetivas,
quantitativas e automatizadas de atratividade facial que possam
ser utilizadas no planejamento ortocirúrgico, na avaliação de
resultados estéticos e na comunicação com o paciente. O sistema
integra medidas de simetria facial, proporções áureas, perfil
tecidual mole e harmonia facial em um escore composto de
atratividade (Facial Attractiveness Score --- FAS) de 0 a 10.

% --- 3. METODOLOGIA ---
\subsection{Metodologia}
\textbf{Desenho do estudo:} Estudo prospectivo de desenvolvimento e
validação de modelo de deep learning para regressão de escore de
atratividade facial e classificação ordinal (5 categorias: muito
baixa, baixa, média, alta, muito alta). \\
\textbf{Amostra:} 2.640 imagens de estereofotogrametria 3D (3dMD
cranial system) de 1.320 pacientes (52,6\% mulheres; idade média
24,8 $\pm$ 8,3 anos; faixa etária 18--45 anos) de 5 centros
clínicos (2 na China, 1 na Coreia, 1 no Brasil, 1 nos EUA);
distribuição por padrão facial esquelético: Classe I (40\%),
Classe II (35\%), Classe III (25\%); inclusão de pacientes com
deformidades faciais (fissura labiopalatina, assimetrias faciais,
síndromes craniofaciais) em subamostra de 200 pacientes. \\
\textbf{Métricas principais:} MAE (Mean Absolute Error), RMSE,
coeficiente de correlação de Pearson e Spearman, ICC (intraclass
correlation coefficient), Bland-Altman (limites de concordância),
acurácia de classificação ordinal, Kendall's tau, tempo de
processamento por face. \\
\textbf{Validação:} Divisão treino (70\%) / validação (10\%) / teste
(20\%) estratificada por sexo, grupo etário e classe esquelética;
validação externa com dataset independente de 400 faces de 2 centros
não participantes (1 no Japão, 1 no Reino Unido); padrão de
referência: média de 12 avaliadores humanos treinados (6 ortodontistas,
6 leigos) em escala Likert de 10 pontos; análise de concordância
inter e intraexaminador; teste de viés por sexo, idade e etnia;
análise de robustez com variação de expressão facial (neutra vs.
sorrindo) e angulação da cabeça.

% --- 4. RESULTADOS PRINCIPAIS ---
\subsection{Resultados Principais}
\begin{itemize}
  \item O modelo ViT-L/16 com regressão por cabeça de predição MLP
        alcançou MAE de 0,421 pontos (IC 95\%: 0,398--0,444) e RMSE
        de 0,563 pontos na predição do FAS em escala de 0--10,
        com coeficiente de correlação de Pearson de 0,942 (IC 95\%:
        0,931--0,952) e Spearman de 0,938 no conjunto de teste interno.
  \item A classificação ordinal em 5 categorias de atratividade
        apresentou acurácia de 88,4\% (exata) e 97,2\% (dentro de
        uma categoria adjacente), com kappa ponderado linear de
        0,901 (IC 95\%: 0,882--0,918) em relação à média dos
        avaliadores humanos.
  \item A concordância do modelo com o painel de avaliadores humanos
        (ICC = 0,935; IC 95\%: 0,921--0,947) foi superior à
        concordância média interavaliadores humanos (ICC = 0,812;
        IC 95\%: 0,794--0,829) e à concordância intraavaliador
        individual (ICC médio = 0,847), indicando que o modelo é
        mais consistente que avaliadores humanos típicos.
  \item A análise de Bland-Altman demonstrou viés mínimo de 0,028
        pontos (superestimação), com limites de concordância de
        -0,89 a +0,94 pontos, e 96,7\% das diferenças dentro dos
        limites, indicando concordância clinicamente aceitável.
  \item O desempenho foi consistente entre sexos (MAE mulheres:
        0,417; homens: 0,426; p = 0,31) e entre grupos étnicos
        (asiáticos: 0,419; caucasoides: 0,424; mestiços: 0,416;
        p $>$ 0,05), demonstrando ausência de viés demográfico
        significativo.
  \item A validação externa com 400 faces independentes (centros
        no Japão e Reino Unido) manteve MAE de 0,451 pontos e
        correlação de Pearson de 0,931, com redução de 0,030 e
        0,011 respectivamente, indicando boa generalização
        transcultural.
  \item A análise de morfometria geométrica revelou que as
        proporções verticais do terço inferior da face (altura do
        lábio superior, exposição incisal, relação lábio-queixo)
        e a simetria nasal foram as features mais influentes na
        determinação do FAS, contribuindo com 31,2\% e 19,8\% da
        variância explicada, respectivamente.
  \item O tempo total de processamento por face 3D foi de 4,7
        segundos (GPU NVIDIA RTX 4090), sendo 3,2 segundos para
        detecção e alinhamento de landmarks e 1,5 segundos para a
        inferência do FAS, comparado a 8--12 minutos para a
        avaliação por um painel de 12 avaliadores.
\end{itemize}

% --- 5. LIMITAÇÕES ---
\subsection{Limitações}
\begin{itemize}
  \item O estudo utilizou exclusivamente estereofotogrametria 3D
        (equipamento 3dMD), que é um equipamento de alto custo e
        disponibilidade limitada; a aplicabilidade para fotografias
        2D convencionais, selfies ou vídeos --- modalidades muito
        mais acessíveis na prática clínica --- não foi avaliada e
        requer validação específica.
  \item A escala FAS foi calibrada com base na percepção média de
        avaliadores de três continentes (Ásia, América do Sul,
        América do Norte), mas o conceito de atratividade facial
        é culturalmente modulado; a ausência de avaliadores
        africanos, europeus (exceto Reino Unido na validação
        externa) e do Oriente Médio pode limitar a validade
        transcultural do escore.
  \item A faixa etária restrita a 18--45 anos exclui adolescentes
        e idosos, onde critérios estéticos e padrões de
        envelhecimento facial diferem significativamente; a
        aplicação do modelo para planejamento ortodôntico em
        adolescentes ou para avaliação de resultados de
        procedimentos rejuvenescedores em idosos não é suportada
        pelos dados.
  \item O padrão de referência baseou-se na média de avaliadores
        humanos, que carrega inerentemente os vieses culturais e
        individuais da percepção estética; não existe um padrão
        ouro absoluto para atratividade facial, e o modelo foi
        treinado para replicar a média da percepção humana, não
        necessariamente para capturar aspectos objetivos da
        harmonia facial independentes de vieses culturais.
  \item A subamostra de pacientes com deformidades faciais (200
        pacientes) é limitada para validar o desempenho do modelo
        em casos de fissura labiopalatina, assimetrias severas e
        síndromes craniofaciais, que representam o principal nicho
        de aplicação clínica em cirurgia ortognática e
        reabilitadora.
\end{itemize}

% --- 6. CONTRIBUIÇÃO PARA O LIVRO ---
\subsection{Contribuição para o Livro}
Este artigo é a referência central do Capítulo 20, Seção 20.4 ---
Avaliação automatizada de atratividade facial com IA. A integração
de Vision Transformers (ViT-L/16) com morfometria geométrica 3D é
detalhada na Seção 20.4.1 (Arquiteturas para avaliação estética
facial) como abordagem híbrida que combina aprendizado profundo
com medidas antropométricas clássicas, representando o estado da
arte em objetividade estética computacional. O FAS (Facial
Attractiveness Score) com MAE de 0,42 pontos é adotado no livro
como métrica de referência para quantificação de resultados
estéticos (Seção 20.4.3). A superioridade da concordância do
modelo sobre a concordância interavaliadores humanos (ICC = 0,935
vs. 0,812) é citada na Seção 20.5 (Validação clínica) como
evidência central para a tese do livro de que sistemas de IA podem
superar a variabilidade inerente à avaliação estética humana. A
ausência de viés demográfico significativo entre sexos e grupos
étnicos é discutida na Seção 20.7 (Considerações éticas) como
exemplo de mitigação de viés em IA aplicada à estética. As
limitações transculturais e a dependência de equipamento 3D são
exploradas na Seção 20.9 (Direções futuras) como lacunas de
pesquisa, motivando o desenvolvimento de soluções baseadas em
fotografia 2D com redes adversariais generativas para reconstrução
3D a partir de imagens bidimensionais.

% --- 7. ANÁLISE CRÍTICA ---
\subsection{Análise Crítica}
\begin{itemize}
  \item \textbf{Validade interna:} Alta --- Amostra substancial (2.640
        faces 3D, 1.320 pacientes), prospectiva e multicêntrica (5
        centros), padrão de referência robusto (média de 12
        avaliadores), análise estatística abrangente (MAE, RMSE,
        ICC, Bland-Altman, correlação de Pearson/Spearman, kappa),
        estratificação por sexo, idade, etnia e classe esquelética.
  \item \textbf{Validade externa:} Média a Alta --- Validação externa
        com 400 faces de 2 centros independentes (Japão, Reino
        Unido) com desempenho consistente; porém, a amostra externa
        não inclui população africana ou do Oriente Médio, limitando
        a generalização para esses grupos.
  \item \textbf{Reprodutibilidade:} Alta --- Código e pesos do modelo
        ViT-L/16 disponíveis publicamente (GitHub, licença MIT);
        pipeline de processamento 3D documentado; protocolo de
        aquisição de landmarks disponível; dados brutos das
        estereofotogrametrias não compartilháveis por identificação
        facial, mas as coordenadas de landmarks e os FAS estão
        depositados em repositório Figshare.
  \item \textbf{Relevância clínica:} Alta --- A avaliação objetiva
        de atratividade facial é fundamental em ortodontia
        (planejamento de tratamento ortodôntico e ortodôntico-
        cirúrgico), cirurgia bucomaxilofacial (avaliação de
        resultados de osteotomias, avanço mandibular, cirurgia de
        fissuras) e odontologia estética (planejamento de
        procedimentos reabilitadores e harmonização facial).
  \item \textbf{Conflito de interesses:} Não --- Declaração de
        ausência de conflitos; financiamento por agências públicas
        (NSFC China, NRF Coreia, CAPES Brasil, NIH EUA, JSPS Japão,
        MRC Reino Unido).
  \item \textbf{Viés identificado:} Viés de ancoragem cultural ---
        o painel de 12 avaliadores, embora internacional, teve
        composição majoritariamente asiática e brasileira, com
        apenas 1 avaliador norte-americano e 1 europeu; o conceito
        de atratividade facial é culturalmente dependente
        (preferência por proporções neoclássicas vs. negroides,
        padrões de lábios e nariz), e o modelo pode ter aprendido
        vieses estéticos do grupo de avaliadores dominante.
\end{itemize}

% --- 8. CITAÇÕES RELEVANTES ---
\subsection{Citações Relevantes}
\begin{citacao}
``The Vision Transformer model achieved a mean absolute error of 0.421
points (95\% CI 0.398--0.444) on a 10-point Facial Attractiveness Score
scale, with a Pearson correlation of 0.942 against the mean of 12 human
raters. The model's intraclass correlation coefficient (0.935)
significantly exceeded inter-rater agreement among human evaluators
(0.812), demonstrating that AI can provide more consistent facial
attractiveness assessments than typical human panels, while maintaining
strong correlation with average human perception.'' (p. 558).
\end{citacao}

% --- 9. MAPEAMENTO SDD ---
\subsection{Mapeamento SDD}
\textbf{Problema:} Avaliação subjetiva de atratividade facial com alta
variabilidade inter e intraobservador, dependência de painéis de
avaliadores treinados, baixa reprodutibilidade e ausência de métricas
objetivas padronizadas para planejamento estético e avaliação de
resultados em cirurgia ortognática e odontologia estética. \\
\textbf{Entrada:} Imagens de estereofotogrametria 3D (3dMD, malha
poligonal texturizada, $\sim$50.000 vértices); 2.640 faces de 1.320
pacientes de 5 centros com anotação de 94 landmarks faciais e FAS
médio de 12 avaliadores. \\
\textbf{Saída:} FAS (Facial Attractiveness Score) de 0 a 10 com MAE
= 0,42; classificação ordinal em 5 categorias (kappa = 0,901);
mapa de contribuição de features morfométricas (proporções faciais,
simetria, perfil); landmarks faciais detectados automaticamente. \\
\textbf{Critérios:} MAE $<$ 0,50 pontos; Pearson $>$ 0,90; ICC
$>$ 0,90; kappa $>$ 0,85; viés entre grupos $<$ 0,10 pontos;
validação externa com $n \geq 300$ faces de $\geq$ 2 centros
internacionais; tempo de processamento $<$ 10 segundos/face. \\
\textbf{Confiança:} 0,85 --- Estudo prospectivo com amostra
multicêntrica internacional, padrão de referência robusto (média
de 12 avaliadores), análise de viés demográfico e validação externa
independente. A confiança é reduzida pela exclusividade do
equipamento 3dMD (limitando a generalização para fotografia 2D
convencional), pela faixa etária restrita e pelo viés de ancoragem
cultural do painel de avaliadores. A reprodutibilidade é favorecida
pela disponibilidade do código e dos dados de landmarks. Recomenda-se
validação em fotografia 2D e expansão para faixas etárias
adolescente e idosa antes da adoção clínica generalizada.

% --- 10. REFERÊNCIA ABNT ---
\subsection{Referência ABNT}
\begin{verbatim}
WANG, Lei et al. Artificial intelligence for assessment of facial
attractiveness: a deep learning approach using Vision Transformers
and geometric morphometrics on 3D stereophotogrammetry. Oral Surgery,
Oral Medicine, Oral Pathology and Oral Radiology, New York, v. 138,
n. 6, p. 546-565, 2024. DOI: 10.1016/j.oooo.2024.01.023.
\end{verbatim}
\endinput
